%! AP Statistics — Unit 3, Lesson 3.6 — Lesson Plan
\documentclass[10pt]{article}
\usepackage{apstats-article}
\usepackage{apstats-boxes}

\newcommand{\UnitNumberName}{Unit 3: Collecting Data \quad}
\newcommand{\LessonNumberName}{Lesson 3.6: Selecting an Experimental Design}

\begin{document}
\begin{center}
    {\Large\bfseries \CourseName: \SchoolYear\ (\MeetingLength) \\
    {\small\bfseries \UnitNumberName \LessonNumberName}}
\end{center}

\begin{tcolorbox}[colback=sky, colframe=navy, boxrule=0.9pt, arc=2mm,
    left=2mm, right=2mm, top=1.5mm, bottom=1.5mm]
    \textbf{Primary Objective:} Students will evaluate the advantages and disadvantages
    of each experimental design and justify which design is most appropriate for a
    given research scenario (Big Idea: VAR).
\end{tcolorbox}

\renewcommand{\arraystretch}{1.6}
\begin{skillbox}[Priority Ideas \& Skills]{goldbox}
\begin{minipage}[t]{0.48\linewidth}
    \textbf{AP Skill 1 --- Selecting Statistical Methods}
    \begin{itemize}
        \item Explain why a particular experimental design is appropriate (VAR-3.D,
              Skill 1.C).
        \item Compare advantages and disadvantages of CRD, block design, and matched
              pairs (VAR-3.D.1).
        \item Match design choice to the research question, resources, and
              nature of the experimental units.
    \end{itemize}
\end{minipage}
\hfill
\begin{minipage}[t]{0.48\linewidth}
    \textbf{Key Understandings}
    \begin{itemize}
        \item CRD: simplest design; best when experimental units are homogeneous.
        \item Block design: reduces variability when a nuisance variable (blocking
              variable) is known; increases precision.
        \item Matched pairs: maximum control of individual differences; useful when
              units can be paired by a relevant characteristic or receive both
              treatments.
        \item Design choice depends on the question, available units, and resources.
    \end{itemize}
\end{minipage}
\end{skillbox}

\begin{skillbox}[{Vocabulary, Concepts, \& Theorems}]{greenbox}
\begin{minipage}[t]{0.48\linewidth}
    \begin{itemize}
        \item \textbf{Completely randomized design (CRD)} --- all experimental units
              randomly assigned to treatments with no blocking
        \item \textbf{Block design} --- units grouped by a blocking variable; treatments
              randomly assigned within each block
        \item \textbf{Blocking variable} --- a variable known to affect the response;
              used to create homogeneous blocks
    \end{itemize}
\end{minipage}
\hfill
\begin{minipage}[t]{0.48\linewidth}
    \begin{itemize}
        \item \textbf{Matched pairs design} --- a randomized block design with blocks of
              size two (pairs matched on a relevant variable, or same unit receives
              both treatments)
        \item \textbf{Nuisance variable} --- a variable that affects the response but is
              not the focus of the experiment; controlled by blocking
        \item \textbf{Precision} --- reduction in variability; blocking increases
              precision by removing the effect of the blocking variable from error
    \end{itemize}
\end{minipage}
\end{skillbox}

\begin{skillbox}[Activate Prior Knowledge \& Spiral Review]{sky}
\begin{minipage}[t]{0.48\linewidth}
    \textbf{Prior Knowledge Check (3--4 min)}
    \begin{itemize}
        \item Recall: ``What are the four principles of a well-designed experiment?''
              (Lesson 3.5)
        \item Recall: ``What is the purpose of blocking?'' (reduce variability due to a
              known nuisance variable)
        \item Transition: ``Today we decide \emph{which} design to use based on the
              research context.''
    \end{itemize}
\end{minipage}
\hfill
\begin{minipage}[t]{0.48\linewidth}
    \textbf{Connections Forward}
    \begin{itemize}
        \item Using experiments to draw causal conclusions $\to$ Lesson 3.7
        \item Scope of inference (random assignment + random selection) $\to$ Lesson 3.7
        \item Two-sample $t$-tests and $z$-tests applied to experiments $\to$ Units 7--8
    \end{itemize}
\end{minipage}
\end{skillbox}

\begin{skillbox}[Hook \& Lesson]{sky}
\begin{minipage}[t]{0.48\linewidth}
    \textbf{Hook (5--7 min)}\\
    Three researchers each propose a different design to test a new pain medication:
    \begin{itemize}
        \item Dr.\ A: randomly assign all 60 patients to drug or placebo.
        \item Dr.\ B: block by gender (30M, 30F), then randomly assign within each.
        \item Dr.\ C: pair each patient with a similar partner; one gets drug, one
              gets placebo.
    \end{itemize}
    Ask: ``Are all three valid? Are they equally good? Which would you choose and
    why?''
\end{minipage}
\hfill
\begin{minipage}[t]{0.48\linewidth}
    \textbf{Lesson Flow (15--18 min)}
    \begin{enumerate}
        \item Review CRD: best when units are homogeneous; simplest to implement.
        \item Block design: use when a nuisance variable is known and measurable;
              increases precision; each block is analyzed separately.
        \item Matched pairs: pairs matched on the most important variable; maximum
              precision; works even with small samples.
        \item Comparison table: CRD vs.\ block vs.\ matched pairs --- key features,
              when to use, pros, cons.
        \item Decision rule: ask (a) Do we know a variable that affects the response?
              (b) Can units be paired? (c) Can same unit receive both treatments?
    \end{enumerate}
\end{minipage}
\end{skillbox}

\begin{skillbox}[Explicit Instruction \& Active Monitoring]{sky}
\begin{minipage}[t]{0.48\linewidth}
    \textbf{Direct Instruction (10 min)}
    \begin{itemize}
        \item Worked example: 40 students are to be assigned to two different test-prep
              programs. Students vary by GPA. Design: (a) CRD, (b) block by GPA
              quartile, (c) matched pairs by GPA.
        \item Ask: if GPA is strongly related to outcome, which design gives the most
              precise estimate of the treatment effect?
    \end{itemize}
\end{minipage}
\hfill
\begin{minipage}[t]{0.48\linewidth}
    \textbf{Active Monitoring}
    \begin{itemize}
        \item Listen for: students recommending matched pairs for every situation
              (it's not always feasible).
        \item Cold-call: ``If the experimental units are all very similar, does blocking
              help? Explain.'' (No; blocking only helps if the blocking variable is
              related to the response.)
    \end{itemize}
\end{minipage}
\end{skillbox}

\begin{skillbox}[Group Work \& Differentiation]{redbox}
\begin{minipage}[t]{0.48\linewidth}
    \textbf{Partner Activity (8--10 min)}\\
    Three research scenarios. Partners:
    \begin{enumerate}
        \item Recommend a design (CRD, block, or matched pairs).
        \item Justify the choice using the context (units, nuisance variable,
              feasibility).
        \item Identify one limitation of the chosen design.
        \item Draw a diagram of the design.
    \end{enumerate}
\end{minipage}
\hfill
\begin{minipage}[t]{0.48\linewidth}
    \textbf{Differentiation}
    \begin{itemize}
        \item \textit{Support}: Decision tree --- ``Is there a known nuisance variable?
              YES: can we pair? YES $\to$ matched pairs. NO $\to$ block design.
              NO $\to$ CRD.''
        \item \textit{Extension}: Justify a design for a multi-site clinical trial
              involving patients from 5 hospitals with different demographics.
        \item \textit{ELL}: Diagrams of each design with color-coded groups.
    \end{itemize}
\end{minipage}
\end{skillbox}

\begin{skillbox}[Individual Work \& Assessment]{redbox}
\begin{minipage}[t]{0.48\linewidth}
    \textbf{Exit Ticket (5 min)}\\
    Scenario: \textit{``A researcher wants to test whether a new fertilizer increases
    tomato yield. She has 24 plots in 4 different fields (6 plots per field). Soil
    quality varies by field.''}
    \begin{enumerate}
        \item What design should she use? Why?
        \item What is the blocking variable?
        \item Draw a diagram of the design showing the 24 plots.
    \end{enumerate}
\end{minipage}
\hfill
\begin{minipage}[t]{0.48\linewidth}
    \textbf{AP-Style Check}\\
    \textit{``A researcher is testing whether a new typing program improves speed.
    She pairs 20 employees by current typing speed, then randomly assigns one in
    each pair to the program and one to a control. This is best described as:''}
    \begin{enumerate}[label=(\Alph*)]
        \item A completely randomized design
        \item A stratified random sample
        \item A matched pairs experiment
        \item A block design with 20 blocks
    \end{enumerate}
    Answer: (C) and (D) are both defensible; on the AP, matched pairs is the
    more specific correct answer.
\end{minipage}
\end{skillbox}

\begin{skillbox}[Reinforcement \& Extension]{goldbox}
\begin{minipage}[t]{0.48\linewidth}
    \textbf{Homework / Practice}
    \begin{itemize}
        \item For 4 research scenarios: recommend a design, justify the choice, and
              draw a diagram.
        \item Reflection: \textit{``When is a CRD preferable to a block design, even
              if you know a variable that affects the response?''}
    \end{itemize}
\end{minipage}
\hfill
\begin{minipage}[t]{0.48\linewidth}
    \textbf{Extension / Preview}
    \begin{itemize}
        \item \textit{Extension}: Research the Salk polio vaccine trial (1954). Explain
              the design choices and why double-blinding was critical.
        \item \textit{Preview}: Lesson 3.7 addresses the logic of inference ---
              how do we know if the difference we observed between treatment groups
              is real or just due to random chance?
    \end{itemize}
\end{minipage}
\end{skillbox}

\end{document}
